\chapter{Resultados}
\label{chap:resultados}

% ==========================================================================
% CAMBIO RESPECTO A 2024
% 2024: primera etapa F=19; IV = +0.716 (1% más multas -> 0.716% más accidentes),
%       significativo al 1%; conclusión: las multas se asocian con MÁS accidentes.
% Ahora: nulo robusto sobre víctimas; el efecto lineal aparente desaparece con
%       Poisson; los fallecidos no se mueven; el "mecanismo" de más enforcement
%       es débil y transitorio.
% ==========================================================================

\section{Tendencias previas}
\label{sec:pretendencias}

La Figura~\ref{fig:pretendencias} contrasta el estudio de evento de los
incidentes C5 y el de las víctimas de tránsito para los tres cambios de política
considerados en la fase de rediseño (febrero de 2021, abril de 2022 y septiembre
de 2023).

\begin{figure}[H]
  \centering
  \includegraphics[width=\textwidth]{509_pretendencias_c5_vs_victimas.jpg}
  \caption{Estudio de evento para incidentes C5 (arriba) y víctimas de tránsito
  (abajo), en los tres cambios de política. \emph{buffer} de 500~m. En cada
  panel se anota el $p$-valor de la prueba conjunta de tendencias previas.}
  \label{fig:pretendencias}
\end{figure}

Con \textbf{incidentes C5} la prueba conjunta se rechaza en los tres eventos y en
todos los \emph{buffers} ($p \approx 10^{-16}$): las colonias cercanas a la
infraestructura vienen en una trayectoria creciente respecto a las lejanas, con
independencia del evento. Este \emph{outcome} no permite una interpretación
causal y se reserva como referencia secundaria.

Con \textbf{víctimas de tránsito} la prueba conjunta no se rechaza en abril de
2022 ni en septiembre de 2023 ($p = 0.11$ y $0.29$ respectivamente). El resto del
análisis se ancla en abril de 2022, el evento con primera etapa plausible
(Sección~\ref{sec:mecanismo}) y con más ventana posterior disponible.

\section{Efecto sobre las víctimas de tránsito}

La Figura~\ref{fig:victimas} presenta el estudio de evento de abril de 2022 para
el total de víctimas y para las víctimas fallecidas, en los tres \emph{buffers}.

\begin{figure}[H]
  \centering
  \includegraphics[width=\textwidth]{510_victimas.jpg}
  \caption{Estudio de evento de la renovación de la Policía de Tránsito
  (abril de 2022). \emph{Outcome}: víctimas de hechos de tránsito.}
  \label{fig:victimas}
\end{figure}

En el \emph{buffer} más estrecho (250~m) se observa un aumento transitorio del
total de víctimas en los primeros meses posteriores al evento, que se desvanece
hacia el mes $+6$. El efecto promedio posterior estimado por mínimos cuadrados
es de $+12$ a $+21\%$ según el \emph{buffer} (Tabla~\ref{tab:did}), marginalmente
significativo. Las \textbf{víctimas fallecidas no muestran ningún efecto} en
ninguna especificación ni horizonte.

\begin{table}[H]
  \centering
  \caption{Efecto promedio posterior (diferencia en diferencias agrupada),
  abril de 2022, ventana $-10/+10$.}
  \label{tab:did}
  \begin{tabular}{llrrr}
  \toprule
  \emph{Outcome} & \emph{buffer} & Efecto (\%) & $p$ & $p$ tend. previas \\
  \midrule
  Víctimas totales & 250 m  & $+21.5$ & 0.001 & 0.39 \\
  Víctimas totales & 500 m  & $+12.6$ & 0.03  & 0.17 \\
  Víctimas totales & 1000 m & $+11.8$ & 0.01  & 0.19 \\
  Víctimas fallecidas & 250 m  & $+14.8$ & 0.42 & 0.17 \\
  Víctimas fallecidas & 500 m  & $+0.7$  & 0.96 & 0.29 \\
  Víctimas fallecidas & 1000 m & $-7.6$  & 0.55 & 0.31 \\
  \bottomrule
  \end{tabular}
\end{table}

\clearpage
\section{Robustez}
\label{sec:robustez}

La Figura~\ref{fig:robustez} recopila el efecto promedio posterior bajo distintas
especificaciones: modelo lineal frente a Poisson de efectos fijos, tratamiento
binario frente a continuo, y ventanas previas alternativas.

\begin{figure}[H]
  \centering
  \includegraphics[width=\textwidth]{603_robustez_att.jpg}
  \caption{Robustez del efecto promedio posterior (abril de 2022). Puntos
  negros: $p<0.05$.}
  \label{fig:robustez}
\end{figure}

El resultado central es que \textbf{el efecto aparente sobre las víctimas totales
no es robusto}. Bajo un modelo de conteo (Poisson con efectos fijos de colonia y
mes), que es la especificación adecuada para un \emph{outcome} con 89\% de
observaciones en cero, el efecto desaparece ($-0.5\%$, $p = 0.93$ con
\emph{buffer} de 500~m; $+8\%$, $p = 0.22$ con 250~m). El efecto positivo que
aparece con mínimos cuadrados está impulsado por unas pocas colonias de vías
primarias con incrementos grandes en términos absolutos. Las víctimas fallecidas
permanecen sin efecto en todas las especificaciones.

\section{Mecanismo: fiscalización diferencial}
\label{sec:mecanismo}

La Figura~\ref{fig:mecanismo} muestra el estudio de evento de las infracciones
alrededor de abril de 2022.

\begin{figure}[H]
  \centering
  \includegraphics[width=\textwidth]{510_mecanismo_multas.jpg}
  \caption{Mecanismo: infracciones cerca de la infraestructura de fiscalización
  alrededor de abril de 2022.}
  \label{fig:mecanismo}
\end{figure}

Se observa un aumento transitorio de las infracciones cerca de la infraestructura
en los primeros meses posteriores a la reforma (del orden de $+15$ a $+35\%$ en
el estudio de evento para $k = 1,\dots,3$, más claro en el \emph{buffer} de
250~m), que regresa a su nivel base hacia el mes $+5$. El efecto promedio
posterior no es estadísticamente significativo ($p \approx 0.3$--$0.6$). La
evidencia de una fiscalización diferencial es, por tanto, \textbf{débil y de
corta duración}.

\section{Réplica en un segundo evento: septiembre de 2023}
\label{sec:replica_sep2023}

Para que la conclusión no dependa de un único cambio de política, se repite el
mismo diseño para la \textbf{reforma para motociclistas de septiembre de
2023}, evento que también cuenta con tendencias previas válidas para las
víctimas de tránsito (Sección~\ref{sec:pretendencias}). Se usa el inventario de
infraestructura de 2022 y la misma ventana $-10/+10$. No se estima el mecanismo
de infracciones: la fuente termina en abril de 2023, antes de este evento, por
lo que ni siquiera el mes de referencia cuenta con cobertura.

\begin{figure}[H]
  \centering
  \includegraphics[width=\textwidth]{511_victimas.jpg}
  \caption{Estudio de evento de la reforma para motociclistas (septiembre de
  2023). \emph{Outcome}: víctimas de hechos de tránsito.}
  \label{fig:victimas_sep2023}
\end{figure}

\begin{figure}[H]
  \centering
  \includegraphics[width=0.85\textwidth]{604_robustez_att.jpg}
  \caption{Robustez del efecto promedio posterior, septiembre de 2023.}
  \label{fig:robustez_sep2023}
\end{figure}

El resultado es consistente con abril de 2022: las víctimas totales no muestran
efecto en ninguna especificación ($-4$ a $+4\%$, nunca significativo). Las
víctimas fallecidas muestran puntos estimados positivos bajo mínimos cuadrados
($+10$ a $+36\%$, marginalmente significativos en un par de especificaciones),
pero \textbf{ninguno sobrevive al modelo de conteo} ($p = 0.43$--$0.52$ con
Poisson). Los incidentes C5 vuelven a mostrar un rechazo fuerte de las
pre-tendencias y no se interpretan.

Que el resultado nulo se repita en dos eventos distintos —con mecanismos
institucionales diferentes y separados por año y medio— hace más creíble que la
ausencia de efecto no es una particularidad de abril de 2022.

\section{Síntesis}

No se encuentra evidencia robusta de que la renovación de la Policía de Tránsito
de abril de 2022 haya reducido las víctimas de hechos de tránsito en las zonas
con infraestructura de fiscalización. Las víctimas fallecidas no cambian en
ninguna especificación; el aumento aparente en víctimas totales no sobrevive a un
modelo de conteo y coincide en momento y duración con el repunte transitorio de
las infracciones, lo que sugiere un efecto de reporte más que un deterioro real
de la seguridad vial. La propia intervención —una fiscalización diferencial
inducida por la reforma— es débil y transitoria, por lo que el resultado nulo
debe leerse como ausencia de evidencia de efecto en un contexto donde el
tratamiento es en sí mismo tenue, más que como un cero estimado con precisión.

La réplica en septiembre de 2023 (Sección~\ref{sec:replica_sep2023}) refuerza
esta lectura: en un evento distinto, con otro mecanismo institucional, el mismo
diseño vuelve a arrojar un nulo robusto una vez que se controla adecuadamente la
dispersión del conteo de víctimas. La ausencia de efecto no parece, entonces,
una particularidad de un solo cambio de política.
